\subsection{指数与对数}

我们用 exp 表示从 C 到 C 的复指数函数（FVR, III, p.~8, 定义 2）。它可微且满足 \(\exp' = \exp\)（FVR, III, p.~9, (26)）；因此它在 C 中全纯。设 A 是一个幺 Banach 代数，\(a\) 是 A 的元素。由 I, p.~65 的命题 2 与 FVR, III, p.~16 的公式 (9)，有

\[
\exp (a) = \sum_ {n = 0} ^ {\infty} \frac {a ^ {n}}{n !}.\tag{9}
\]

由于 \(\|a^{n}\|\leqslant\|a\|^{n}\)，我们看出 \(\|\exp(a)\|\leqslant\exp(\|a\|)\)，并且级数 (9) 在 A 的每个球中一致收敛。映射 \(a\mapsto\exp(a)\)（从 A 到 A）是全纯的（I, p.~76 的命题 10）。我们也时常用 \(e^{a}\) 表示 \(a\in A\) 的指数。

当 \(a\) 是 Banach 空间 E 的自同态时，像这样定义的指数 \(\exp(a)\)——在 Banach 代数 \(\mathcal{L}(\mathrm{E})\) 中——与 FVR, IV, p.~27, 定义 1 中定义的那个相一致，依据同前命题 7 (3)。

对 A 的每个与 a 可交换的元素 b，我们还有

\[
\exp (a + b) = \sum_ {n = 0} ^ {\infty} \frac {b ^ {n}}{n !} \exp (a),
\]

（I, p.~77 的命题 11），由此得

\[
\exp (a + b) = \exp (a) \cdot \exp (b).\tag{10}
\]

特别地，\(\exp(a)\) 可逆，且

\[
\exp (a) ^ {- 1} = \exp (- a).\tag{11}
\]

设 B 是由这样的 \(z \in C\) 组成的集合：\(-\pi < Jz < \pi\)。设 F 是 C 中区间 \(R_{-}\) 的余集。指数到 B 上的限制通过过渡到子空间诱导出从 B 到 F 上的一个双射（FVR, III, p.~10, n°7），其逆双射将记作 log。

若 \(a \in A\) 满足 \(\mathrm{Sp}_{\mathrm{A}}(a) \subset \mathrm{F}\)，就可以作出 A 的元素 \(\log(a)\)。我们有 \(\mathrm{Sp}_{\mathrm{A}}(\log(a)) \subset \mathrm{B}\)，并且

\[
\exp (\log (a)) = a\tag{12}
\]

（由 I, p.~75 的命题 8）。反之，设 \(b\) 是 A 中满足 \(\mathrm{Sp}_{\mathrm{A}}(b) \subset \mathrm{B}\) 的元素。我们有 \(\mathrm{Sp}_{\mathrm{A}}(\exp(b)) \subset \mathrm{F}\)，并且

\[
\log (\exp (b)) = b\tag{13}
\]

（同前）。

特别地，若 \(a \in A\) 满足 \(\varrho(a) < 1\)，则 \(\mathrm{Sp}_{\mathrm{A}}(1 - a) \subset \mathrm{F}\)，我们可以作出 \(\log(1 - a)\)。对 \(n \geqslant 1\)，第 \(n\) 阶导数（映射 \(z \mapsto \log(1 - z)\) 的）是 \(z \mapsto -(n - 1)! (1 - z)^{-n}\)。因此 \(z \mapsto \log(1 - z)\) 在点 0 处的幂级数展开为

\[
\log (1 - z) = - \sum_ {n = 1} ^ {\infty} \frac {z ^ {n}}{n},
\]

对 \(|z| < 1\) 成立（VAR, R1, p.~30, 3.3.9）。由 I, p.~65 的命题 2，我们得到

\[
\log (1 - a) = - \sum_ {n = 1} ^ {\infty} \frac {a ^ {n}}{n}.\tag{14}
\]

命题 12. --- 设 A 是一个交换幺 Banach 代数。指数映射的像是 A 的可逆元群 G 中含单位元的连通分支。

公式 (10) 与 (11) 证明 \(\exp(A)\) 是 G 的一个子群。由前文（见公式 (12)），这个子群包含以 1 为中心、半径为 1 的开球。因此它是 G 的开子群，从而是闭子群。此外，A 是连通的，且映射 \(a \mapsto \exp(a)\) 连续，所以 \(\exp(A)\) 连通。于是 \(\exp(A)\) 是 G 的含单位元的连通分支。

